%% ma: L11-B17-0016
%% lop: 11
%% chuong: 5
%% bai: 17
%% dang: 11.17.3
%% loai: TL
%% muc: TH
%% dap_an: "-"
%% nguon: "tai-lieu-hoc-tap-mon-toan-lop-11-theo-chuong-trinh-giao-duc-pho-thong-2018.pdf (Lớp Toán Thầy Duy, Phạm Lê Duy, THPT Chu Văn An), Chương 3 Bài 3, Dạng 3, Ví dụ 16 (tr.393-398)"
%% kien_thuc: "Lớp 11 Bài 17 (hàm số liên tục)"
%% trang_thai: chinh-thuc
%% ghi_chu: "Đã duyệt 03/11/2026."
\begin{ex}
Chứng minh rằng phương trình $x^3+(m+3)x^2+(1-m)x-1=0$ luôn có nghiệm với mọi $m$.
\loigiai{\sloppy Đặt $f(x)=x^3+(m+3)x^2+(1-m)x-1$, liên tục trên $\mathbb R$. $f(0)=-1<0$ và $f(1)=1+m+3+1-m-1=4>0$ với mọi $m$ nên $f(0)f(1)<0$. Vậy phương trình luôn có nghiệm thuộc $(0;1)$.}
\end{ex}
